\chapter{Supplementary Results: Experiment D}
\label{appendix:FullWTResults}
% Referenced from: Sec 5.2.3 ("See Appendix ... for the raw PSDs of each case
% plotted on their own axes, and full tables of raw results"), and the captions
% of Tables 5.2 and 5.3 ("Full raw statistics are given in ...").

\section{Raw PSDs per Wind-tunnel Setpoint}
\label{appendix:WTPSDs}
\begin{figure}[H]
	\centering
	\includegraphics[width=0.7\linewidth]{"../Data Acquisition/figures/psd_windtunnel_cases"}
	\caption{Grid of PSD plots for each wind-tunnel setpoint}
	\label{fig:psdwindtunnelcases}
\end{figure}

\section{Raw Noise and Peak Statistics}
\label{appendix:WTRawStats}
% >>> generated: WTRawStats -- written by "Data Acquisition/make_appendixG_tables.py"; edits between these markers are overwritten
\autoref{tab:WTRawStats-fog} and \autoref{tab:WTRawStats-mist} give the un-normalised counterparts of \autoref{tab:D_snr} and \autoref{tab:D_mist_avg}: the same Doppler peaks and noise bands, but with the statistics taken on the raw PSD rather than on the PSD divided by its running-median floor. Because the raw floor slopes across the noise band, its standard deviation here includes that slope as well as the random fluctuation, which makes the raw mean $+\,5\sigma$ threshold the more conservative of the two.

\begin{table}[H]
	\centering
	\small
	\begin{tabular}{@{}S[table-format=2.1] S[table-format=2.3] S[table-format=1.2] S[table-format=1.3] S[table-format=1.2] S[table-format=1.2] S[table-format=1.2] S[table-format=2.2]@{}}
		\toprule
		{Setpoint} & {$f_D$} & \multicolumn{4}{c}{PSD ($10^{-12}$\,V$^2$/Hz)} & {Peak/mean} & {Excess} \\
		\cmidrule(lr){3-6}
		{(m/s)} & {(MHz)} & {floor mean} & {floor $\sigma$} & {mean $+\,5\sigma$} & {at $f_D$} & {(dB)} & {($10^{-7}$\,V$^2$)} \\
		\midrule
		3.1 & 3.637 & 3.01 & 0.078 & 3.41 & 4.73 & 1.96 & 10.25 \\
		5.0 & 5.868 & 3.45 & 0.075 & 3.82 & 4.79 & 1.43 & 8.87 \\
		6.0 & 6.892 & 2.44 & 0.031 & 2.59 & 3.78 & 1.90 & 7.55 \\
		7.0 & 7.997 & 2.30 & 0.040 & 2.50 & 3.49 & 1.82 & 7.04 \\
		8.0 & 9.280 & 4.30 & 0.082 & 4.71 & 5.36 & 0.95 & 5.99 \\
		9.0 & 10.611 & 2.74 & 0.042 & 2.95 & 3.60 & 1.19 & 6.08 \\
		10.0 & 11.664 & 4.20 & 0.054 & 4.46 & 5.18 & 0.92 & 6.00 \\
		11.8 & 13.837 & 3.94 & 0.053 & 4.20 & 4.51 & 0.59 & 5.27 \\
		21.0 & 24.674 & 3.69 & 0.052 & 3.95 & 4.22 & 0.59 & 4.01 \\
		\bottomrule
	\end{tabular}
	\caption{D: raw (un-normalised) noise and Doppler-peak statistics for dense fog-seeded air}
	\label{tab:WTRawStats-fog}
\end{table}

\begin{table}[H]
	\centering
	\small
	\setlength{\tabcolsep}{4.5pt}
	\begin{tabular}{@{}S[table-format=1.1] S[table-format=1.0] S[table-format=2.3] S[table-format=1.2] S[table-format=1.3] S[table-format=1.2] S[table-format=1.2] S[table-format=1.2] S[table-format=2.2]@{}}
		\toprule
		{Setpoint} & {Captures} & {$f_D$} & \multicolumn{4}{c}{PSD ($10^{-12}$\,V$^2$/Hz)} & {Peak/mean} & {Excess} \\
		\cmidrule(lr){4-7}
		{(m/s)} & {averaged} & {(MHz)} & {floor mean} & {floor $\sigma$} & {mean $+\,5\sigma$} & {at $f_D$} & {(dB)} & {($10^{-7}$\,V$^2$)} \\
		\midrule
		5.0 & 1 & 5.730 & 4.48 & 0.090 & 4.93 & {4.80\textsuperscript{\dag}} & 0.30 & 4.38 \\
		5.0 & 2 & 5.716 & 3.83 & 0.087 & 4.27 & {4.08\textsuperscript{\dag}} & 0.28 & 4.45 \\
		5.0 & 3 & 5.726 & 4.00 & 0.081 & 4.41 & {4.29\textsuperscript{\dag}} & 0.30 & 3.75 \\
		\bottomrule
	\end{tabular}

	\vspace{2pt}
	{\raggedright\footnotesize \textsuperscript{\dag}\,PSD at $f_D$ below the raw mean $+\,5\sigma$ threshold. The detection verdicts in \autoref{tab:D_mist_avg} use the floor-normalised PSD instead, on which the \SI{5}{m/s} peak clears its threshold with three captures.\par}
	\caption{D: raw (un-normalised) noise and Doppler-peak statistics for light mist-seeded air.}
	\label{tab:WTRawStats-mist}
\end{table}


\section{Velocity Accuracy: Peak Centroid vs Setpoint}
\label{appendix:velocityTable}

\autoref{tab:WTVelocity} gives the data behind \autoref{fig:figvelocityboth}: for each setpoint, the Doppler-peak centroid, the line-of-sight velocity it implies, the velocity expected from the setpoint and the \SI{35}{\degree} beam angle, and the error between the two.

\begin{table}[H]
	\centering
	\small
	\begin{tabular}{@{}S[table-format=2.1] S[table-format=2.3] S[table-format=2.2] S[table-format=2.2] S[table-format=+2.1, retain-explicit-plus]@{}}
		\toprule
		{Setpoint $U$} & {$f_D$} & {Measured $v_{\mathrm{LOS}}$} & {Expected $v_{\mathrm{LOS}}$} & {Error} \\
		{(m/s)} & {(MHz)} & {$= f_D\lambda/2$ (m/s)} & {$= U\cos 35^\circ$ (m/s)} & {(\%)} \\
		\midrule
		\multicolumn{5}{@{}l}{\textit{Collimator facing downstream (original rig)}} \\
		3.1 & 3.637 & 2.82 & 2.54 & +11.0 \\
		5.0 & 5.868 & 4.55 & 4.10 & +11.0 \\
		6.0 & 6.892 & 5.34 & 4.91 & +8.7 \\
		7.0 & 7.997 & 6.20 & 5.73 & +8.1 \\
		8.0 & 9.280 & 7.19 & 6.55 & +9.8 \\
		9.0 & 10.611 & 8.22 & 7.37 & +11.6 \\
		10.0 & 11.664 & 9.04 & 8.19 & +10.4 \\
		11.8 & 13.837 & 10.72 & 9.67 & +10.9 \\
		21.0 & 24.674 & 19.12 & 17.20 & +11.2 \\
		\multicolumn{5}{@{}l}{\quad Fit through origin: \SI{1.170 \pm 0.003}{MHz} per m/s, +10.7\,\% above the expected \SI{1.057}{MHz} per m/s} \\
		\midrule
		\multicolumn{5}{@{}l}{\textit{Collimator facing upstream (rotated rig)}} \\
		3.9 & 4.269 & 3.31 & 3.19 & +3.6 \\
		5.3 & 5.733 & 4.44 & 4.34 & +2.3 \\
		6.1 & 6.802 & 5.27 & 5.00 & +5.5 \\
		6.3 & 6.816 & 5.28 & 5.16 & +2.4 \\
		7.5 & 8.422 & 6.53 & 6.14 & +6.2 \\
		8.0 & 9.105 & 7.06 & 6.55 & +7.7 \\
		9.0 & 10.265 & 7.95 & 7.37 & +7.9 \\
		10.0 & 11.337 & 8.79 & 8.19 & +7.3 \\
		21.0 & 23.864 & 18.49 & 17.20 & +7.5 \\
		\multicolumn{5}{@{}l}{\quad Fit through origin: \SI{1.130 \pm 0.006}{MHz} per m/s, +6.9\,\% above the expected \SI{1.057}{MHz} per m/s} \\
		\bottomrule
	\end{tabular}
	
	\caption{D: Doppler-peak centroid and line-of-sight velocity against the wind-tunnel setpoint for both rig orientations, in dense fog-seeded air}
	\label{tab:WTVelocity}
\end{table}

